% !TEX root = ../Projektdokumentation.tex
\section{Einleitung}
\label{sec:Einleitung}


\subsection{Projektumfeld} 
\label{sec:Projektumfeld}
\begin{itemize}
	\item JournawayGmbH ist ein 2017 gegründetes Reiseunternehmen und bietet Reise-komplett-Pakete an. Derzeit werden rund 160 Mitarbeiter beschäfitgt.
	\item Das Projekt dient primär dazu die Mitarbeiter im Kundensupport zu entlasten und eine Source of Truth zu etablieren.
\end{itemize}


\subsection{Projektziel} 
\label{sec:Projektziel}
\begin{itemize}
	\item Bei dem Projekt geht es darum einen .NET Microservice zu entwickeln, der mithilfe von RabbitMQ Kundendaten zweier Systeme synchronisiert.
	\item Am Ende soll erreicht werden, dass die Kundendaten auf beiden Systemen identisch sind.
\end{itemize}


\subsection{Projektbegründung} 
\label{sec:Projektbegruendung}
\begin{itemize}
	\item Das Projekt ist sinnvoll, da dadurch veraltete Informationen zu Kunden und manuelle Datenpflege vermieden werden können.
	\item Dadurch das der Kundenlogin neu hinzugefügt wird ist es notwendig, dass die Kundendaten auf allen Systemen korrekt dargestellt werden.
\end{itemize}


\subsection{Projektschnittstellen} 
\label{sec:Projektschnittstellen}
\begin{itemize}
	\item Die Anwendung interagiert über Event Queues durch RabbitMQ mit einem anderen Microservice 'user-service'. Außerdem kommuniziert die Anwendung über einen HttpClient mit dem CRM 'HubSpot'. Schließlich wird mit einer PostgreSql Datenbank interagiert in der die Synchronisationshistorien gespeichert werden.
	\item Das Projekt wird innerhalb des Customer Retention Teams bearbeitet und besprochen. 
	\item Dadurch, dass es sich um einen autonomen Synchronisationsservice handelt gibt es keine direkten Benutzer. Die Synchronisation wird entweder bei der Eingabe von Kundendaten auf der Webseite -oder HubSpot durch Kunden oder Mitarbeiter ausgelöst.
	\item Das Ergebnis wird am Ende in einem Sprintreview vor anderen Mitarbeitern präsentiert
\end{itemize}


\subsection{Projektabgrenzung} 
\label{sec:Projektabgrenzung}
\begin{itemize}
	\item Die Aufnahme der Anwendung in den Livebetrieb ist nicht Teil des Projektes.
\end{itemize}
